\section{\texorpdfstring{\(C = 6ND\)}{C = 6ND}, và cái nó bỏ quên}
\label{sec:ch09-sau-nd}

\index{quy luật scaling!công thức 6ND}%
Hai bài của mục~\ref{sec:ch09-hai-quy-luat} đều đặt mọi kết luận của mình lên
một xấp xỉ, và nó là một dòng: huấn luyện một mô hình \(N\) tham số trên \(D\)
token tốn chừng \(6ND\) phép tính dấu chấm động.

Chỗ con số 6 ra, Kaplan và cộng sự viết ở mục 2.1:

\begin{quote}
\enquote{Accounting for the backwards pass (approximately twice the compute as
the forwards pass), we then define the estimated non-embedding compute as
\(C \approx 6N\) floating point operations per training
token.}~\autocite{kaplan2020scaling}
\end{quote}

Mỗi tham số dự đúng một phép nhân-cộng cho mỗi token, và một phép nhân-cộng là
2 FLOP, nên lượt xuôi tốn \(2N\). Lượt ngược tốn gấp đôi, \(4N\). Cộng lại là
\(6N\) trên mỗi token, nhân \(D\) token.

Chương~\ref{ch:chi-phi-va-song-song} đã dựng một phép đếm FLOP chính xác cho
một tầng và kiểm nó bằng bộ đếm toán tử của chính torch. Nên câu hỏi ở đây trả
lời được: \(6ND\) sai bao nhiêu, và sai ở đâu.

\subsection*{Bảng 1 của Kaplan đếm một trong hai tích bậc hai}

\index{quy luật scaling!bảng 1 thiếu một tích}%
Trước khi đo sai số, có một chuyện về chính bảng ấy.

Attention có \emph{hai} phép nhân ma trận mang \(n^2\), và
mục~\ref{sec:ch08-dem} đã đếm cả hai: \(\mat{Q}\mat{K}^\top\) để ra
\tn{ma trận điểm số}{score matrix}, rồi phép lấy tổ hợp \(\mat{V}\) theo trọng
số softmax. Mỗi phép tốn
\(2n^2 d\), tổng \(4n^2 d\) một tầng.

Bảng 1 của Kaplan có đúng một hàng mang \(n_{\mathrm{ctx}}\), tên là
\code{Attention: Mask}, và hàng ấy ghi \(2 n_{\mathrm{layer}}
n_{\mathrm{ctx}} d_{\mathrm{attn}}\). Hàng \code{Total} của bảng ghi
\(C_{\mathrm{forward}} = 2N + 2 n_{\mathrm{layer}} n_{\mathrm{ctx}}
d_{\mathrm{attn}}\), và phương trình 2.2 trong thân bài mang đúng cùng hệ số
ấy. Bảng và công thức nhất quán với nhau, và cùng đếm một tích.

\begin{measured}{bảng 1 của Kaplan đặt cạnh phép đếm chính xác, mỗi token một tầng}
\begin{minted}[fontsize=\footnotesize]{text}
d_model  n_ctx   kaplan attn   exact attn    ratio  kaplan total  exact total
768      512     786,432       1,572,864     2.00   14,942,208    15,728,640
1600     1024    3,276,800     6,553,600     2.00   64,716,800    67,993,600
12288    2048    50,331,648    100,663,296   2.00   3,674,210,304 3,724,541,952
\end{minted}

Cấu hình chuẩn của chính bài, \(d_{\mathrm{attn}} = \dff/4 = \dmodel\).

Đo bằng \code{experiments/ch09\_counts.py} tại tag \code{ch09}.
\end{measured}

Tỷ lệ đúng 2.00 ở cả ba cấu hình. Nửa tham số của hai phép đếm thì khớp nhau
chính xác, nên chỗ lệch không phải một cách quy ước khác về FLOP mà đúng là một
tích bị thiếu.

\textbf{Ai đúng thì không phải chuyện đọc kỹ hơn.} \code{tests/test\_ch09.py}
đem phép đếm dạng đóng so với \code{FlopCounterMode}, tức bản ghi của torch về
các toán tử thật sự đã chạy, trên một stack thật ở ba kích thước; chênh lệch là
0 ở cả ba. Đây là cùng phép kiểm chéo mục~\ref{sec:ch08-dem} dùng cho một tầng,
nâng lên cho cả chồng. Phụ lục F của Hoffmann và cộng sự cũng liệt kê cả hai
tích, từng dòng một, dưới hai cái tên \code{Key @ Query logits} và
\code{Softmax @ query reductions}~\autocite{hoffmann2022chinchilla}.

Và điều đáng nói là chuyện ấy không quan trọng. Số hạng bị đếm thiếu chính là số
hạng mà \(6N\) vứt đi hoàn toàn ở bước sau. Bài nói thẳng ra là nó vứt, cùng
mục 2.1: \enquote{Since we primarily study models where \(d_{\mathrm{model}} \gg
n_{\mathrm{ctx}}/12\), we do not include context-dependent terms in our training
compute estimate.} Một hằng số 2 trong một số hạng sắp bị bỏ thì không đi tới
đâu cả. Tôi in nó ra vì bảng 1 là chỗ bài viết phép đếm ra thành công thức, nên
nó là chỗ một người dựng phép đếm FLOP của mình sẽ chép từ đó, và người ấy sẽ
mang theo đúng chỗ thiếu này.

\subsection*{Sai số của 6ND là đại lượng chương trước đã đo}

\index{quy luật scaling!sai số của 6ND}%
Câu hỏi thật thì ở chỗ khác: bỏ số hạng ấy đi thì mất bao nhiêu.

Viết ra thì gọn. Mỗi token, mỗi tầng, bốn phép chiếu tốn \(8d^2\), mạng truyền
thẳng tốn \(4d\dff\), và hai tích attention tốn \(4nd\). Ở \(\dff = 4d\) thì
hai số hạng đầu gộp thành \(24d^2\), nên phần attention chiếm

\begin{equation}
\label{eq:ch09-phan-attention}
\frac{4nd}{24d^2 + 4nd} = \frac{n}{6d + n}.
\end{equation}

Số tầng biến mất, vì mọi số hạng đều nhân với nó. Còn lại một hàm của đúng tỷ
số \(n/d\), và nó bằng một nửa tại \(n = 6d\).

\(6d\) là con số của mục~\ref{sec:ch08-dem}. Ở đó \(n = 2d + \dff\) là chỗ phần
bậc hai vượt nửa một tầng, và ở \(\dff = 4d\) thì \(2d + 4d = 6d\). Hai chương
tới cùng một ngưỡng từ hai phía: chương trước hỏi khi nào phần bậc hai đắt bằng
phần còn lại, chương này hỏi khi nào xấp xỉ của Kaplan bỏ sót một nửa. Cùng một
\(n\).

Từ \eqref{eq:ch09-phan-attention}, tỷ lệ giữa phép đếm chính xác và \(2N\) là
xấp xỉ \(1 + n/(6d)\), và \(D\) triệt tiêu khỏi nó vì cả hai vế đều tuyến tính
theo \(D\).

\begin{measured}{6ND so với phép đếm chính xác, ở bốn cấu hình đã công bố}
\begin{minted}[fontsize=\footnotesize]{text}
configuration                   n_ctx   ratio    1 + n/(6d)  attn share
GPT-1        L12 d768           512     1.1095   1.1111      0.1000
GPT-2 1542M  L48 d1600          1024    1.1059   1.1067      0.0964
GPT-3 175B   L96 d12288         2048    1.0277   1.0278      0.0270
GPT-3 175B, 4x context          8192    1.1110   1.1111      0.1000
GPT-3 175B, 16x context         32768   1.4443   1.4444      0.3077
\end{minted}

Đo bằng \code{experiments/ch09\_counts.py} tại tag \code{ch09}.
\end{measured}

Ba chuyện đọc ra từ bảng ấy.

Thứ nhất, xấp xỉ \(1 + n/(6d)\) bám số đo trong vòng 0.002 ở khắp bảng, nên
\eqref{eq:ch09-phan-attention} không phải một ước lượng thô mà là chính đại
lượng ấy, sai lệch chỉ vì \(N\) còn mang bias và layer norm.

Thứ hai, \(6ND\) luôn thấp hơn thực tế, không bao giờ cao hơn. Nó bỏ đi một số
hạng dương và không bỏ gì khác.

Thứ ba, và đây là chỗ đáng dừng: sai số ở GPT-3 là 2.8\% còn ở GPT-1 là 11\%.
Không phải vì GPT-3 lớn hơn. Là vì GPT-3 \emph{rộng} hơn nhiều so với độ dài
ngữ cảnh của nó. Cùng mô hình GPT-3 ấy, giữ nguyên mọi tham số và chỉ kéo ngữ
cảnh lên 32768, sai số thành 44\%. Đại lượng quyết định là \(n/d\), đúng như
\eqref{eq:ch09-phan-attention} nói, và không phải quy mô.

Đó cũng là lý do các bài \tn{quy luật scaling}{scaling law} dùng được \(6ND\) mà
vẫn tin được:
chúng khớp trên các mô hình mà \(n/d\) nhỏ. Điều kiện Kaplan viết ra,
\(d_{\mathrm{model}} \gg n_{\mathrm{ctx}}/12\), là đúng đại lượng ấy nói theo
cách khác. Một chương sau, ở chỗ ngữ cảnh dài, điều kiện ấy hỏng, và
chương~\ref{ch:vi-tri-va-ngu-canh} là chương phải sống với chỗ hỏng ấy.

\subsection*{Vì sao \texorpdfstring{\(N\)}{N} trừ embedding ra}

\index{quy luật scaling!N không tính embedding}%
Còn một chi tiết định nghĩa mà bỏ qua thì mọi con số dưới đây lệch. Kaplan viết
ở mục 1.3: \(N\) là \enquote{the number of model parameters, excluding all
vocabulary and positional embeddings}, và mục 2.1 nói lý do:
\enquote{we will see that this produces significantly cleaner scaling laws}.

Nghe như chuyện sổ sách. Không phải.

\begin{measured}{bảng embedding chiếm bao nhiêu phần của mô hình}
\begin{minted}[fontsize=\footnotesize]{text}
model        total               embedding       non-embedding    emb share
GPT-1        116,536,320         31,480,320      85,056,000       0.2701
GPT-2 1542M  1,557,611,200       82,049,600      1,475,561,600    0.0527
BERT-base    109,482,240         23,837,184      85,645,056       0.2177
GPT-3 175B   174,604,259,328     642,723,840     173,961,535,488  0.0037
\end{minted}

Hàng GPT-1 lấy từ điển 40,478 của checkpoint phát hành, vì bài không in kích
thước từ điển và mục~\ref{sec:ch09-cung-mot-khoi} đã nói chỗ đó. Chọn đầu kia
của dải, 40,000, thì phần embedding thành 0.2678 thay vì 0.2701, nên chỗ không
chắc ấy không đụng gì tới điều bảng này nói.

Đo bằng \code{experiments/ch09\_counts.py} tại tag \code{ch09}.
\end{measured}

Bảng embedding là 27\% của GPT-1 và 0.37\% của GPT-3. Hệ số 73, trên đúng dải
quy mô mà các quy luật được khớp. Khối stack lớn theo \(d^2\) còn bảng embedding
lớn theo \(d\), nên hai đại lượng ấy tách nhau ra càng lúc càng xa.

Hệ quả thực hành thì gọn: một quy luật khớp theo \(N\) không tính embedding mà
đem áp cho tổng số tham số là đang áp qua một hệ số trôi 73 lần. Và
Hoffmann và cộng sự đếm \emph{ngược lại} với Kaplan, ghi ở phụ lục F:
\enquote{we also count embeddings matrices in the total parameter count}. Hai
bài, hai định nghĩa của cùng một chữ \(N\), và
mục~\ref{sec:ch09-hai-quy-luat} là chỗ chuyện ấy có hậu quả.
